\documentclass{article}
\usepackage[T1]{fontenc}
\usepackage{lmodern}
\usepackage{graphicx}
\usepackage{amsmath,amssymb}
\usepackage[a4paper,margin=2cm]{geometry}
\usepackage[absolute,overlay]{textpos}
\setlength{\TPHorizModule}{1mm}
\setlength{\TPVertModule}{1mm}
\usepackage[indonesian]{babel}
\usepackage{hyperref}
\setlength{\emergencystretch}{2em}
% unit: Halaman 22

\begin{document}

% === Halaman 22 (llm) ===

\section*{Perkalian Titik}

\begin{itemize}
    \item Perkalian titik: \(kP = Q\)
    
    Ket: \(k\) adalah skalar, \(P\) dan \(Q\) adalah titik pada kurva eliptik
    
    \item Perkalian titik diperoleh dengan perulangan dua operasi dasar kurva eliptik yang sudah dijelaskan:
    \begin{enumerate}
        \item Penjumlahan titik \((P + Q = R)\)
        \item Penggandaan titik \((2P = R)\)
    \end{enumerate}
    
    \item Contoh: \(k = 3 \rightarrow 3P = P + P + P\) atau \(3P = 2P + P\)
    
    \(k = 23 \rightarrow kP = 23P = 2(2(2(2P) + P) + P) + P\)
\end{itemize}

\end{document}
